\documentclass[a4paper,12pt]{article}


%\usepackage{amssymb}
\usepackage{fontspec}
%有這個套件才能自由設定字體

\setmainfont[BoldFont=cwTeXQHei-Bold]{cwTeX Q Ming}
%cwTeX 預設無粗體，所以必須特別為英文字體設定粗體字體
%若不會用到 \textbf 指令可以去掉[BoldFont=cwTeXQHei-Bold] 

\newfontfamily\mysm{CMU Typewriter Text}
%打程式碼的時候會另外用到 type writer font
%取用時打 {\mysm YOURCODE} 即可
%沒裝 CMU 自行可以用 Courier New 來取代（較細）


\XeTeXlinebreaklocale "zh"
\XeTeXlinebreakskip = 0pt plus 1pt
%使用 XeLaTeX 的主要斷行設定

\setlength{\parskip}{0.5cm}
\linespread{1.15}
%根據個人愛好設定行距與段落間距

\begin{document}
\author{r99245002~~~~林致翰}
\title{
The Transistor (invented in 1947) and MOSFETs - Physics - Historical development - Present status and future perspectives}

\maketitle

在電路學上電晶體屬於三阜的元件，正由於它比起一般只有兩隻接腳的電阻、電容或者電感之類的被動元件多了一個控制端的緣故，我們可以利用電晶體做出一些控制的效果，比如控制某兩端點的電流來決定另外兩端點間是否導通，做出類似開關的效果，又或者加上適當的偏壓以及限流電阻作成放大電路來做出放大電路。我們甚至可以結合多個電晶體組成邏輯電路來處理一些基本數位信號，利用電晶體的非線性效果來處理判別高電位以及低電位信號，換言之可以用來處理 0 或者 1，加上振盪器製造時脈來觸發邏輯閘的信號（當然振盪器本身也是由電晶體構成的正回受電路），最後形成電腦運算單元的雛型。在電晶體以前是屬於真空管的時代，電晶體的功能真空管當然也有，但兩者的體積、信號響應速度、耗電與使用壽命不可相提並論，我們甚至可以大膽的說沒有電晶體這項發明就不可能發展出現代微型高速高效率的電腦與各式各樣輕巧的電子產品，台灣足以自傲的半導體產業技術源頭最終還是要追溯到如何改良電晶體使其更有效率。

電晶體有相當多不同的架構，當然功能也有所不同，比如 BJT （Bipolar Junction Transistor， 其中的 Bipolar Junction 指的是電晶體中 PNP 或者 NPN 結構所具有的 PN 介面），或者數位積體電路常用的 MOSFET （metal–oxide–semiconductor field-effect transistor, 金氧半導體場效電晶體，其中的金屬作為導電極，現代半導體製程中則由矽層取代金屬層）MOSFET 利用 Oxide 層隔開導電控制極與低濃度摻雜的半導體通道，當你在這個控制極，一般稱為閘級施加電壓的時候會吸引半導體通道的載子進而改變這個通道運蘇載子的能力，隨著這個通道兩端的壓差呈現不同導電特性，所以稱為場效電晶體）。由於只需要在閘級上施加電壓，不需要像 BJT 一樣施加電流就能控制另外兩個電極間的電性，所以 MOSFET 操作功耗較 BJT 來的小，所以現代數位積體電路大多使用 MOSFET 構成，至於類比、通訊電路則考量到操作頻段（換句話說與載子傳輸，訊號，響應速度有關）則可能由 BJT 或者兩者混合特性互補的 BiCMOS 來製作。除了一般邏輯電路外，MOSFET 只要在閘級加上儲存電荷的電容結構就能組成記憶單元，比如說需要隨時補充電荷才能維持記憶狀態的動態記憶體 (DRAM)，又或者是電容結構以絕緣層加工成浮動電極，使得電荷在閘級不施加偏壓時，儲存在閘級的電荷可以長期保存，不會洩漏到載子通道之中的快閃記憶體 （Flash Memory）。

現代 MOSFET 的製作通常是使用矽作為基板（也有使用三五族元素，但是產量與市場份額較矽相比則遜色不少），使用離子部植的方式在基板上摻雜雜質使結構出現多於的電子或電洞作為載子，至於電極幾何形狀與層狀結構則可透過蝕刻的方式完成。想要提高電路密度，將 MOSFET 微型化的關鍵在於閘極的寬度，換言之就是閘極控制的半導體通道長度，這個長度越短，你能在單位面積下置入 MOSFET 的數量就越多，一般常聽到的 xx nm 製程(比如說台積電的 28 nm 製程， intel cpu 從 Sandy Bridge 的 32 nm 製程到最新 Ivy Bridge 的 22 nm 製程) 指的是半導體工藝下閘級可以作到的最小線寬，一般來說這個長度局限於蝕刻流程中光罩曝光的繞射極限，除了使用傳統光罩曝光來蝕刻出電極形狀，目前也有使用多電子束直接加工基板的技術（考量到效率問題只使用一個電子束是不行的，需要成千上萬束才足夠），只是成本可能較 EUV 曝光蝕刻的方式高，但是考慮到曝光光源的波長一但縮小到某個程度時幾乎沒有適合的反射光學元件可以使用（或者說極為昂貴），多電子束的蝕刻技術不失為一個有商業價值的好選擇。

另外一個 MOSFET 需要考慮的問題便是信號響應速度以及功耗，這個就與電極與基板材料有關，比如說使用 high k 材料作為絕緣層來改善絕緣層厚度，降低洩漏電流來取得更好的功耗表現，改變基板材質或者以二次摻雜的方式來提高半導體通道的載子遷移率，這些相關的技術基本需要具備固態物理領域知識的人去解決，所以修習這門導論課程的人未來如果想從事相關產業可以從這方面入手，如果想從光罩蝕刻降低線寬方面切入半導體產業則最好多接觸一點 EUV 光學或者電子束技術比較好。我認為未來 MOS 的發展必須兩個方面同時下手才有辦法接近原子等級線寬的技術，偏重任一方面都沒有辦法作到，好的光源與好的材料必須同時具備才有辦法做出足夠優秀的電晶體。

不過我認為 MOSFET 的發展也快到頭了，雖然目前還有 3D IC 等新技術來解決電晶體蜜度問題，但是高電流密度帶來的散熱問題與製程良率都是我們會面臨到的困境，也許下一代的商用高速運算晶片必須改用純光學方式來解決（silicon optics），但現階段看來積體化的效益還是以半導體電信號處理方式佔優勢，光晶片只在傳輸速度上佔優勢，若不能積體化則未來幾個 decade 還是傳統半導體的天下。

\end{document}